\section{Discussion}
\label{sec:discussion}

\subsection{Interpretation of the results}

The matched experiments show that zero-shot forecasting can be competitive
with a practical supervised baseline on retail data. Chronos-2 is the most
consistent model in our panels, and its advantage survives several changes to
the M5 sample and LightGBM specification. This is useful evidence against the
view that a foundation model is relevant only when no trained baseline is
available. At the same time, the differences are generally modest, the panels
are small, and other foundation models do not share the same ranking. The
result is best understood as a reason to include a strong zero-shot model in a
retail evaluation, not as a reason to replace established systems without
testing them.

The public benchmarks support the same qualified interpretation. Foundation
models compare particularly well on Scaled Quantile Loss, suggesting value for
applications that require a predictive distribution. Their advantage is much
smaller on point-forecast measures in fev-bench. This difference may reflect
the objectives used in pre-training, the ability of the models to represent
uncertainty, or the composition of the tasks. The available two-suite sample
cannot distinguish among these explanations.

Model size is not a useful guide within the range we tested. Chronos-2 is
stronger than the smaller Chronos-Bolt-Tiny, but Moirai and TiRex often compare
well with the larger TimesFM. Nor do the results reveal a common horizon
effect. Checkpoint selection based only on parameter count or an overall
leaderboard would therefore miss differences that appear at the dataset and
metric level.

\subsection{Information available to the models}
\label{subsec:data-leakage}\label{subsec:covariate-gap}

Two information asymmetries complicate the comparison. First, foundation
models are trained on large collections that may include public retail data.
Chronos documentation distinguishes in-domain and zero-shot evaluations, and
GIFT-Eval was designed to reduce leakage, but public model cards do not provide
complete dataset-level pre-training manifests. We cannot exclude exposure to
M5 or Favorita. Rohlik~v2 is newer and we found no documented exposure in the
model papers, although absence from the documentation is not proof that the
data were unseen.

Second, our LightGBM models use retail covariates while the foundation models
are univariate. This gives the supervised baseline access to prices,
promotions, and calendar information, but it also reflects a realistic use
case. The close results for several models on Favorita, where promotions and
holidays are important, suggest that a covariate-aware foundation model would
be a valuable next comparison. The present experiments cannot show whether it
would improve on the univariate forecasts.

These asymmetries work in different directions: pre-training exposure may
favour the foundation model, whereas target-specific covariates may favour
LightGBM. They cannot be corrected by a statistical adjustment to the
reported errors. More transparent pre-training records and matched
covariate-aware experiments are needed.

\subsection{Accuracy measures and intermittent demand}
\label{subsec:perseries-vs-aggregate}

The M5 results demonstrate the importance of aggregation. Mean per-series
\WAPE{} gives equal influence to every item, including series with almost no
demand in the test period. A small absolute miss can then produce a very large
ratio. Aggregate \WAPE{} weights observations by demand volume and is less
sensitive to this problem. \WRMSSE{} uses historical variation and economic
weights and answers a different question again.

None of these measures is universally preferable. Item availability may
justify giving low-volume series substantial weight, while aggregate planning
may reasonably focus on volume. The implication is that an evaluation should
state the decision level and report enough information to show whether its
ranking is driven by a sparse tail. In datasets with intermittent demand, a
volume-weighted measure and a per-series or stratified analysis provide a more
complete picture than either alone.

\subsection{Limitations and further work}
\label{subsec:limitations}\label{sec:limitations}

The structured search was screened by one reviewer and does not meet the
standard of a duplicate-screened systematic review. Its coverage is also
shaped by what benchmark authors make extractable. The 498 comparisons come
from only two suites, and multiple rows within a suite share tasks, baselines,
and actual values. For this reason, we report distributions of results rather
than a pooled inferential estimate. The prediction-level fev-bench rerun
improves uncertainty estimation for \WAPE{} against Seasonal Naive, but not
for the other metrics, GIFT-Eval, or comparisons with the strongest tree
baseline.

The experiments cover three datasets and 100-series panels. Favorita is
represented by high-volume series in the main panel, and the more intermittent
check is limited to M5. The LightGBM models are credible practical baselines
but not competition-winning, retailer-specific pipelines. A 200-trial search
was completed only for the main M5 panel, and uncertainty intervals condition
on the selected baseline. Larger panels and high-budget checks on Rohlik and
Favorita would show whether the observed differences persist.

All foundation models are used zero-shot and no checkpoint exceeds 200 million
parameters. The study therefore says nothing about fine-tuning, the largest
models, or the use of exogenous variables by models that support them. Local
inference was measured on a MacBook Air rather than a production GPU system.
Runtime and energy observations cannot be extrapolated to large batches or
service deployments.

Retail coverage is also incomplete. Fashion, fresh produce, e-commerce, and
spare-parts demand have life cycles and intermittency patterns not fully
represented by M5, Rohlik, and Favorita. Future work should prioritise larger
matched panels in these settings, prediction-level benchmark releases,
covariate-aware foundation models, and a common hardware protocol. These
extensions would be more informative than adding further unmatched
leaderboard results.

The main lessons may nevertheless extend to other demand applications:
comparisons should use the same observations, a baseline with a credible
information set, and a measure tied to the decision. Whether the numerical
advantages transfer beyond retail remains an empirical question.
\label{subsec:generalisability}
